\section{State preparation and gate noise}
\label{sec:problem}

The objective is to approximate a known normalized target
$\ket{t}\in\mathbb C^{2^n}$ starting from $\ket{0^n}$.
For a circuit with $L$ elementary gates, we use the ordering
\begin{equation}
 C(\boldsymbol\theta)=G_L(\theta_L)\cdots G_1(\theta_1).
 \label{eq:circuit}
\end{equation}
Each trainable elementary gate has one scalar parameter. Fixed gates can
also occur in the comparator circuits. The noiseless fidelity and loss are
\begin{equation}
 F_0(\boldsymbol\theta)=
 |\bra{t}C(\boldsymbol\theta)\ket{0^n}|^2,
 \qquad \mathcal L_0=1-F_0.
 \label{eq:ideal_fidelity}
\end{equation}
All target amplitudes are available to the simulator. This access assumption
also permits exact evaluation of gate-removal damage.

After each elementary gate, a Pauli channel acts on its support. The channel
is parameter independent. For window $i$, let $P_{i,a}$ denote its nonidentity Pauli
outcomes and let $\lambda_{i,a}\geq0$ be their rates. With
$\kappa_i=\sum_a\lambda_{i,a}$, define
\begin{equation}
 p_{i,0}=e^{-\kappa_i},\quad
 p_{i,a}=(1-e^{-\kappa_i})\frac{\lambda_{i,a}}{\kappa_i}.
 \label{eq:jump_law}
\end{equation}
\begin{equation}
 \mathcal E_i(\rho)=p_{i,0}\rho+
 \sum_a p_{i,a}P_{i,a}\rho P_{i,a}^{\dagger}.
 \label{eq:channel}
\end{equation}
For $\kappa_i=0$, the channel is the identity. Each window contains either
no jump or one listed jump, and outcomes in different windows are sampled
independently. This defines the complete channel used here. It is not the
full repeated-jump evolution generated by a Lindblad equation.

For the profile used in the main comparison, each acted-on site has $X$,
$Y$, and $Z$ rates of $10^{-4}$ each. A two-qubit gate additionally has
correlated $XX$ and $ZZ$ rates of $1.5\times10^{-3}$ each. Thus
$\kappa_i=3\times10^{-4}$ for a one-qubit window and
$\kappa_i=3.6\times10^{-3}$ for a two-qubit window.
These coefficients are simulation inputs, not average gate-error
probabilities or a current device calibration. The historical profile name
is IBM. Idle noise and readout noise are absent.

Starting from $\rho_0=\ket{0^n}\bra{0^n}$, exact propagation gives
\begin{equation}
 \rho_i=\mathcal E_i\!\left(G_i\rho_{i-1}G_i^{\dagger}\right),
 \qquad F=\bra{t}\rho_L\ket{t}.
 \label{eq:noisy_fidelity}
\end{equation}
In the logical-gate study, a noise window follows every original elementary
operation, including a rotation whose angle is zero. Deleting a gate removes
both the operation and its noise window. Parameter changes alone do not
remove windows. Logical $R_{XX}$, $R_{YY}$, and $R_{ZZ}$ rotations and logical
CX gates all receive a two-qubit window. This convention fixes a simulation
model; their physical decomposition costs can differ.

The separate native study instead places the same arity-dependent channel
after operations in a compiled $\{R_Z,\mathrm{SX},X,\mathrm{CX}\}$ schedule.
The schedule is frozen during refinement. The two studies therefore answer
different implementation questions. Neither is a measurement on a quantum
processor.

Training propagates dense state vectors along sampled histories where the
method requires trajectories. Exact feedback methods and final assessment
use density matrices. At six qubits these objects contain 64 amplitudes and
$64\times64$ entries, respectively. The random-MPS target label describes
target generation, not the simulation backend.
